\documentclass[conference]{IEEEtran}

\usepackage{cite}
\usepackage{amsmath,amssymb}
\usepackage{booktabs}
\usepackage{graphicx}
\usepackage{xcolor}
\usepackage{url}
\usepackage{microtype}
\usepackage{multirow}
\usepackage{array}
\usepackage{balance}

\begin{document}

\title{Assurance-Aware Design Space Exploration\\
for Zero-Trust Avionics SoC Security}

\author{\IEEEauthorblockN{[Author names omitted for double-blind review]}}

\maketitle

\begin{abstract}
Modern avionics platforms increasingly consolidate safety-critical and
security-sensitive functions onto shared System-on-Chip (SoC) platforms.
This consolidation improves size, weight, and power characteristics, but it
also creates new attack surfaces because a compromised IP core or bus master
can reach other assets through shared interconnects. This paper presents an
assurance-aware design space exploration (DSE) workflow for avionics SoCs
that jointly synthesizes security feature assignments and zero-trust
enforcement policies under hardware constraints. The workflow uses Answer Set
Programming (ASP) to assign per-component security and logging features
subject to resource, power, and latency limits, then synthesizes firewall and
policy-server placement together with least-privilege access control. The
method exposes excess privileges, trust-anchor gaps, and topology-driven
resilience weaknesses early in design. Evaluation on a primary 8-component
SoC drone case and a secondary 11-component UAV translation shows that the
workflow identifies actionable architectural findings while quantifying
tradeoffs across security and implementation cost. On the primary case, the
zero-trust overlay reduces baseline risk to~51 with 2 firewalls and 1 policy
server, while exposing 9 excess privileges and 6 no-root-of-trust gaps. On
the UAV case, incumbent-bounded continuation proves the \textit{min\_resources}
point at 8,410 LUTs and residual risk~39, while the topology still reveals a
single-string internal-bus weakness and 22 excess privileges.
\end{abstract}

\begin{IEEEkeywords}
Zero-Trust Architecture, System-on-Chip Security, Design Space Exploration,
Answer Set Programming, Avionics
\end{IEEEkeywords}

\section{Introduction}

Avionics architectures are moving toward deeper hardware consolidation.
Processors, accelerators, DMA engines, sensors, and mission functions that
were once separated across boards are increasingly realized on shared
System-on-Chip (SoC) platforms. That transition improves SWaP
characteristics, but it also changes the cybersecurity problem. A compromise
on one bus-connected component can create architectural reachability to many
other assets, including safety-critical ones, through the shared interconnect
fabric.

This creates a design-time problem rather than a purely operational one.
Avionics architects must choose which hardware security features to deploy,
where to place enforcement, and how to preserve latency and resource
constraints. Manual hardening patterns help, but they do not explore the full
design space and do not prove whether a selected architecture is
resource-feasible or close to optimal. That gap is especially relevant to
zero-trust avionics discussions, which increasingly emphasize explicit
verification and least privilege but rarely address on-chip synthesis under
hardware constraints~\cite{nist_zta,ztas_dasc}.

This paper presents an assurance-aware DSE workflow for zero-trust avionics
SoC security. Given a single architectural model, the workflow produces three
practical outputs:
\begin{enumerate}
\item security and logging feature assignments under resource, power, and
latency constraints;
\item zero-trust enforcement placement and least-privilege policy synthesis;
\item resilience-oriented architectural findings derived from compromise and
failure scenarios.
\end{enumerate}

The paper focuses on the design-time core: feature assignment and policy
synthesis. Runtime adaptation and certification-evidence generation remain
important, but they are outside the scope of this DASC submission.

The main contributions are:
\begin{itemize}
\item an ASP-based formulation for avionics SoC security feature selection
under hardware constraints;
\item automated zero-trust policy synthesis that places firewalls and policy
servers while exposing excess privilege and trust gaps;
\item evaluation on two avionics-relevant SoC topologies showing how the
workflow surfaces topology-driven weaknesses that manual review can miss.
\end{itemize}

\section{Related Work}

\textbf{SoC hardware security.}
Hardware-enforced isolation on SoCs is well established through
TrustZone-style partitioning~\cite{arm_tz}, FPGA isolation flows~\cite{xilinx_idf},
and secure interconnect mechanisms~\cite{fiorin_noc}. These mechanisms are
important building blocks, but they do not solve the synthesis problem of
deciding which protections to deploy and where to place them under
implementation constraints.

\textbf{Avionics security and ZTA.}
Avionics security is shaped by DO-326A~\cite{do326a} process concerns and
partitioning assumptions such as ARINC~653~\cite{arinc653}. Zero-trust
principles are defined by NIST~SP~800-207~\cite{nist_zta}. Hasan
et~al.~\cite{hasan_saecps,hasan_jsys} applied zero-trust design patterns
manually to cyber-physical architectures and highlighted the need for tool
support at design time.

\textbf{Risk and exact design exploration.}
NIST~SP~800-30~\cite{nist_sp80030} motivates structured risk reasoning, while
Aven~\cite{aven_riskmatrix} warns against over-interpreting multiplicative
risk matrices on ordinal scales. We therefore use an additive residual-risk
model for standalone components. For exact combinational design-space search,
ASP provides solver-backed optimality and strong logical expressiveness for
reachability and policy reasoning~\cite{gebser_asp}.

Our position is narrower than a general avionics assurance framework. The
contribution here is a constraint-based workflow that jointly reasons about
feature assignment, zero-trust enforcement placement, and hardware budgets at
the SoC architecture level.

\section{System Model and Workflow}

We model the SoC as a directed graph $G=(V,E)$ whose vertices represent bus
masters, receiver IPs, buses, firewalls, and policy servers. Edges represent
interconnect links. Each receiver component carries:
\begin{itemize}
\item asset impact values for read and write operations;
\item candidate security and logging features;
\item per-feature cost in LUTs, FFs, BRAM, and power;
\item allowable read and write latency;
\item trust-domain and trust-anchor information.
\end{itemize}

The workflow has two synthesis phases and one analysis phase.

\begin{figure}[t]
\centering
\fbox{\parbox{0.95\columnwidth}{
\centering
\textbf{Workflow Overview}\\[3pt]
Architectural model $\rightarrow$ Phase~1 feature assignment $\rightarrow$
Phase~2 zero-trust synthesis $\rightarrow$ Phase~3 resilience analysis}}
\caption{Workflow used in the DASC submission version. The paper focuses on
design-time synthesis and compact resilience analysis, not runtime adaptation.}
\label{fig:workflow}
\end{figure}

\subsection{Phase~1: Security Feature Assignment}

Each component receives exactly one security feature and one logging feature.
Table~\ref{tab:features} summarizes the feature set used in the experiments.

\begin{table}[t]
\caption{Feature Summary Used in Phase~1}
\label{tab:features}
\centering
\scriptsize
\begin{tabular}{l l r r r r r}
\toprule
Feature & Type & Lat. & LUTs & FFs & BRAM & mW \\
\midrule
zero\_trust       & sec & 3 & 1200 & 850 & 0 & 24 \\
dynamic\_mac      & sec & 6 & 950  & 680 & 0 & 18 \\
mac               & sec & 4 & 650  & 420 & 0 & 12 \\
zero\_trust\_logger & log & 2 & 520  & 480 & 2 & 11 \\
some\_logging     & log & 1 & 180  & 220 & 1 & 4 \\
no\_logging       & log & 1 & 0    & 0   & 0 & 0 \\
\bottomrule
\end{tabular}
\end{table}

For standalone components, the residual-risk model is additive:
\begin{equation}
R = \max(0,\; \mathit{impact} + \mathit{domain\_bonus}
- \mathit{security\_protect} - \mathit{log\_protect})
\label{eq:risk}
\end{equation}
For redundant groups, the encoding uses a normalized group-compromise model so
that common-mode dependence remains visible even when individual members
appear well protected.

The phase-1 solver minimizes total residual risk subject to:
\begin{itemize}
\item one security feature and one logging feature per component;
\item LUT, FF, DSP, LUTRAM, BRAM, and power budgets;
\item per-asset latency limits;
\item effective per-asset risk caps.
\end{itemize}

The tool supports three strategy modes: \textit{max\_security},
\textit{min\_resources}, and \textit{balanced}.

\subsection{Phase~2: Zero-Trust Policy Synthesis}

Phase~2 injects the phase-1 design as background facts and synthesizes
firewall placement, policy-server placement, least-privilege access policies,
and trust-anchor gap findings. The key constraint is architectural
mediation: if a lower-trust master can topologically reach a protected
target, the path must be mediated by an enforcement point. The solver also
compares topology-implied access against declared \texttt{access\_need}
facts, flagging excess privileges and missing privileges.

\subsection{Phase~3: Resilience Analysis}

Phase~3 is used here only for compact architectural assessment. It executes
compromise and failure scenarios against the synthesized design and reports
total scenario risk, worst-case amplification relative to baseline, service
degradation or loss, and control-plane concentration effects. This DASC
version keeps resilience analysis as a design-time diagnostic rather than a
full runtime response framework.

\section{Primary Case Study: SoCDrone (TC9)}

The primary case is an 8-component avionics-style SoC with 2 masters, 2
buses, a 5-member redundancy group, 3 standalone IPs, and one
latency-sensitive safety-critical component.

\begin{figure}[t]
\centering
\fbox{\parbox{0.95\columnwidth}{
\centering
\textbf{TC9 Topology}\\[3pt]
Two buses; a 5-member redundancy group on noc0; standalone IPs on noc1;
synthesized firewalls guarding each domain.}}
\caption{Primary TC9 case study. The shared bus that supports the redundancy
group also creates a common-mode weakness.}
\label{fig:tc9}
\end{figure}

\subsection{Security Feature Results}

Table~\ref{tab:tc9} summarizes the measured three-strategy results.

\begin{table}[t]
\caption{TC9 Strategy Comparison}
\label{tab:tc9}
\centering
\scriptsize
\begin{tabular}{lrrr}
\toprule
Metric & max\_sec & min\_res & balanced \\
\midrule
Total risk         & 51    & 69    & 51 \\
LUTs used          & 9,340 & 5,380 & 6,840 \\
FFs used           & 7,100 & 3,580 & 4,880 \\
Power (mW)         & 184   & 100   & 132 \\
Firewalls placed   & 2     & 2     & 2 \\
Policy servers     & 1     & 1     & 1 \\
Excess privileges  & 9     & 9     & 9 \\
No-RoT trust gaps  & 6     & 6     & 6 \\
Worst amplification& 3.0$\times$ & 2.839$\times$ & 3.0$\times$ \\
Runtime (s)        & 6.366 & 10.185 & 19.250 \\
\bottomrule
\end{tabular}
\end{table}

These results show a clear security-footprint tradeoff. \textit{min\_resources}
reduces LUT use from 9,340 to 5,380, but raises total risk from 51 to 69.
\textit{balanced} recovers the same total risk as \textit{max\_security}
while cutting LUT use to 6,840.

\subsection{Zero-Trust Synthesis Findings}

Phase~2 produces stable structural findings across all strategies. The solver
places 2 firewalls and 1 policy server, and it consistently identifies 9
excess privileges, 6 no-root-of-trust gaps, and a concentrated policy-control
dependency. The most important result is that the DMA master is materially
over-privileged relative to declared need. That finding is topological, not
merely a consequence of weak local protection settings.

\subsection{Resilience Findings}

The worst scenario for TC9 is compromise of the full redundancy group, which
produces a 3.0$\times$ amplification over the secured baseline for
\textit{max\_security} and \textit{balanced}. This exposes an avionics lesson:
the 5-member group improves availability only if the interconnect itself
remains healthy. Because the members share the same bus, the architecture
still contains a common-mode weakness. The case also shows why design-time
co-optimization matters: the safety-critical path is latency constrained,
which prevents arbitrarily strong protection choices from being assigned
without violating timing.

\section{Secondary Portability Case: DARPA/UAV Translation}

To show portability, the workflow was also applied to an 11-component UAV
architecture translated from the DARPA CASE surveillance system used in prior
manual hardening work~\cite{hasan_saecps,hasan_jsys,darpa_case}. The
translated SoC contains 3 masters, 4 buses, no redundancy groups, 4
safety-critical components, and one bridge from an untrusted ingress path
into the trusted internal bus.

\begin{table}[t]
\caption{DARPA/UAV Portability Summary}
\label{tab:uav}
\centering
\scriptsize
\begin{tabular}{lrrr}
\toprule
Metric & max\_sec & min\_res & balanced \\
\midrule
Total risk         & 0      & 39     & 0 \\
LUTs used          & 18,170 & 8,410  & 16,530 \\
FFs used           & 14,120 & 6,160  & 12,640 \\
Power (mW)         & 367    & 160    & 333 \\
Firewalls placed   & 1      & 1      & 1 \\
Policy servers     & 1      & 1      & 1 \\
Excess privileges  & 22     & 22     & 22 \\
No-RoT trust gaps  & 10     & 10     & 10 \\
Worst amplification& N/A    & 2.062$\times$ & N/A \\
Runtime (s)        & 12.228 & 3,788.100 & 12.482 \\
Optimality proven  & yes    & yes    & yes \\
\bottomrule
\end{tabular}
\end{table}

This case complements TC9. Under \textit{max\_security}, the architecture is
strong enough to drive total residual risk to zero, but the topology still
reveals 22 excess privileges and a narrow internal-bus trust boundary. The
incumbent-bounded continuation run ultimately proved the
\textit{min\_resources} point at 8,410 LUTs, 6,160 FFs, 160~mW, and residual
risk~39. The final assignment is almost the cheapest feasible pattern:
\texttt{mac} everywhere, \texttt{some\_logging} only where required, and
\texttt{no\_logging} where the hard constraints allow it. Even so, the
architecture still exposes a 2.062$\times$ worst-case amplification under
policy-server compromise. In other words, even when feature assignment is
minimized successfully, the bus organization can still dominate the security
story.

This section should be read as portability evidence rather than as a strict
``tool versus Hasan'' comparison. The manual AADL-based work and the present
SoC DSE operate at different abstraction levels. The value of the current
tool is that it exposes privilege and topology findings automatically once
the SoC interconnect model is given.

\section{Discussion and Limitations}

Two results matter most for avionics architects.

First, zero-trust feature assignment alone is not enough. The workflow
repeatedly exposes structural weaknesses such as over-privileged DMA access,
policy-control concentration, and narrow trust boundaries at bus bridges.
These are design-time architecture issues, not just misconfigured
access-control lists.

Second, the tool provides a disciplined alternative to manual hardening trade
studies. Instead of choosing protections qualitatively, the architect can ask
whether a given protection mix is feasible under power, area, and latency
constraints, and whether the resulting topology still contains avoidable
privilege or trust gaps.

The present paper has three main limits. It does not include a head-to-head
baseline against CP-SAT or heuristic search. It models architecture-level
timing through latency caps rather than cycle-accurate execution analysis.
And it stops at design-time synthesis rather than full certification evidence
or runtime adaptation. A separate 36-instance synthetic sweep is encouraging
but still places the current tractability boundary between medium and large
families: all 9 small cases solved exactly, all 9 medium cases produced
solutions, but only 1 of 9 large cases and none of the 9 extra-large cases
solved under the current timeout budget. Those are important next steps, but
they do not reduce the value of the current design-time workflow.

\section{Conclusion}

This paper presented an assurance-aware DSE workflow for zero-trust avionics
SoC security. The workflow jointly synthesizes security feature assignments
and enforcement placement under hardware constraints, then uses resilience
scenarios to expose topology-driven weaknesses. The primary TC9 case showed a
clear tradeoff between residual risk and FPGA footprint while revealing 9
excess privileges and 6 no-root-of-trust gaps. The secondary DARPA/UAV case
showed that even when the minimum-resource point is pushed down to a proven
8,410 LUT design, the architecture still retains serious privilege and
trust-boundary weaknesses. The separate synthetic sweep further showed exact
solutions on all small instances and mixed exact/best-found behavior on
medium ones, with a clear transition to failure between the medium and large
families under the current encoding and timeout budget.

The central conclusion is that avionics SoC security should be treated as a
constrained architecture-synthesis problem, not a manual hardening exercise.
ASP provides a practical way to solve that problem while keeping resource
budgets, timing limits, and zero-trust enforcement visible to the designer.

\balance
\bibliographystyle{IEEEtran}
\bibliography{references}

\end{document}
